\documentclass[UTF8,fontset=fandol,openany,zihao=-4]{ctexbook}
% Shared lecture-note preamble for Big Data and Management Decision-Making.
% Chapter files follow latex_config.md and remain independently compilable.

\usepackage[a4paper,margin=2.6cm]{geometry}
\usepackage{amsmath,amssymb,amsthm,mathtools,bm}
\usepackage{xcolor}
\usepackage[most]{tcolorbox}
\usepackage{booktabs,longtable,array}
\usepackage{enumitem}
\usepackage{hyperref}
\usepackage{listings}
\usepackage{caption}
\usepackage{tikz}
\usetikzlibrary{arrows.meta,positioning,shapes.geometric}

\newcommand{\BDMFandolPath}{/usr/local/texlive/2025/texmf-dist/fonts/opentype/public/fandol/}
\setCJKmainfont[
  Path=\BDMFandolPath,
  UprightFont=FandolSong-Regular.otf,
  BoldFont=FandolSong-Regular.otf,
  ItalicFont=FandolKai-Regular.otf,
  BoldItalicFont=FandolKai-Regular.otf
]{FandolSong-Regular.otf}
\setCJKsansfont[
  Path=\BDMFandolPath,
  UprightFont=FandolSong-Regular.otf,
  BoldFont=FandolSong-Regular.otf
]{FandolSong-Regular.otf}
\setCJKmonofont[
  Path=\BDMFandolPath,
  UprightFont=FandolSong-Regular.otf,
  BoldFont=FandolSong-Regular.otf
]{FandolSong-Regular.otf}
\setmonofont{Menlo}
\linespread{1.13}

\hypersetup{
  colorlinks=true,
  linkcolor=blue!60!black,
  citecolor=blue!60!black,
  urlcolor=blue!60!black,
  pdfauthor={大数据与管理决策基础},
  pdfsubject={大数据与管理决策基础课程讲义}
}

\ctexset{
  chapter = {
    format = \huge\mdseries,
    name = {第,讲},
    number = \arabic{chapter},
    aftername = \quad
  },
  section = {format = \Large\mdseries},
  subsection = {format = \large\mdseries},
  subsubsection = {format = \normalsize\mdseries}
}

\setlist[itemize]{leftmargin=2em,itemsep=0.25em,topsep=0.25em}
\setlist[enumerate]{leftmargin=2em,itemsep=0.25em,topsep=0.25em}

\definecolor{BDMBlue}{RGB}{22,44,73}
\definecolor{BDMTeal}{RGB}{22,119,119}
\definecolor{BDMGray}{RGB}{76,84,95}
\definecolor{BDMLight}{RGB}{245,247,250}
\definecolor{BDMAmber}{RGB}{181,111,35}
\definecolor{CodeBg}{RGB}{248,248,248}

\lstdefinestyle{pythonstyle}{
  basicstyle=\ttfamily\small,
  backgroundcolor=\color{CodeBg},
  keywordstyle=\color{BDMBlue},
  commentstyle=\color{BDMGray},
  stringstyle=\color{BDMAmber},
  showstringspaces=false,
  breaklines=true,
  frame=single,
  rulecolor=\color{black!20},
  columns=fullflexible
}

\lstdefinestyle{sqlstyle}{
  language=SQL,
  basicstyle=\ttfamily\small,
  backgroundcolor=\color{CodeBg},
  keywordstyle=\color{BDMBlue},
  commentstyle=\color{BDMGray},
  stringstyle=\color{BDMAmber},
  showstringspaces=false,
  breaklines=true,
  frame=single,
  rulecolor=\color{black!20},
  columns=fullflexible
}

\lstset{style=pythonstyle}
\renewcommand{\lstlistingname}{代码}
\renewcommand{\bibname}{参考文献}

\theoremstyle{definition}
\newtheorem{definition}{定义}[chapter]
\newtheorem{example}[definition]{例}
\newtheorem{exercise}[definition]{习题}
\newtheorem{convention}[definition]{约定}

\theoremstyle{remark}
\newtheorem{remark}[definition]{评注}

\theoremstyle{plain}
\newtheorem{theorem}[definition]{定理}
\newtheorem{proposition}[definition]{命题}
\newtheorem{lemma}[definition]{引理}
\newtheorem{corollary}[definition]{推论}

\tcolorboxenvironment{definition}{
  colback=blue!2!white,
  colframe=blue!45!black,
  boxrule=0.4pt,
  arc=1.2mm,
  breakable
}

\tcolorboxenvironment{theorem}{
  colback=orange!3!white,
  colframe=orange!55!black,
  boxrule=0.4pt,
  arc=1.2mm,
  breakable
}

\tcolorboxenvironment{proposition}{
  colback=orange!2!white,
  colframe=orange!50!black,
  boxrule=0.4pt,
  arc=1.2mm,
  breakable
}

\tcolorboxenvironment{lemma}{
  colback=orange!2!white,
  colframe=orange!45!black,
  boxrule=0.4pt,
  arc=1.2mm,
  breakable
}

\tcolorboxenvironment{corollary}{
  colback=orange!2!white,
  colframe=orange!45!black,
  boxrule=0.4pt,
  arc=1.2mm,
  breakable
}

\newtcolorbox{chapterbox}{
  colback=gray!5!white,
  colframe=gray!55!black,
  boxrule=0.4pt,
  arc=1.2mm,
  breakable
}

\newtcolorbox{modernbox}[1][应用接口]{
  colback=green!3!white,
  colframe=green!45!black,
  boxrule=0.4pt,
  arc=1.2mm,
  title={#1},
  fonttitle=\mdseries,
  breakable
}

\newcommand{\R}{\mathbb R}
\newcommand{\N}{\mathbb N}
\newcommand{\E}{\mathbb E}
\newcommand{\Prob}{\mathbb P}
\newcommand{\dd}{\,\mathrm d}
\newcommand{\eps}{\varepsilon}
\newcommand{\abs}[1]{\left\lvert #1\right\rvert}
\newcommand{\norm}[1]{\left\lVert #1\right\rVert}
\newcommand{\argmin}{\operatorname*{arg\,min}}
\newcommand{\argmax}{\operatorname*{arg\,max}}


\title{《大数据与管理决策基础》\\[0.5em]\large 第3讲：SQL、KPI 与可审计指标层}
\author{课程讲义}
\date{}

\lstset{style=sqlstyle}

\begin{document}

\maketitle
\tableofcontents

\setcounter{chapter}{2}
\chapter{SQL、KPI 与可审计指标层}
\label{ch:week03-sql-kpi}

\begin{chapterbox}
本讲讨论如何把企业数据模型转化为可查询、可解释、可复核的管理指标。SQL 是关系数据之上的声明式查询语言，但在管理决策场景中，它不只是“取数工具”，而是连接业务对象、事实粒度、维度属性、时间窗口、分子分母和指标责任的表达语言。本讲先解释 SQL 查询的关系语义，再讨论连接、过滤、聚合、条件聚合和窗口函数如何形成销售额、毛利率、准时交付率、延期天数等 KPI；随后进一步说明比率型指标、事实表粒度、指标字典和可审计 SQL 的最低要求。课程目标不是让学生背诵 SQL 片段，而是训练学生写出能够进入管理汇报、数据产品和审计流程的指标查询。
\end{chapterbox}

\begin{modernbox}[课程接口]
本讲把第二讲的数据模型推进到指标层。后续数据质量、统计看板、预测模型、因果评估和优化决策都需要以稳定指标为输入。若 SQL 指标口径不稳定，后续模型即使复杂，也只是在不可靠的数字上继续叠加复杂性。
\end{modernbox}

\section{学习目标}

第一讲建立了从经营症状到行动闭环的主线，第二讲说明企业运营数据如何被建模为实体、事件、状态、事实表和维度表。第三讲进入更具体的问题：当数据模型已经存在时，如何用 SQL 从这些表中构造管理指标？

完成本讲后，学生应能够：
\begin{enumerate}
  \item 解释 SQL 查询与关系模型中选择、投影、连接和聚合之间的关系；
  \item 区分事务粒度、分析粒度和指标展示粒度；
  \item 使用连接、分组聚合和条件聚合计算销售额、毛利、准时交付率和平均延期天数；
  \item 正确处理比率型指标的分子和分母，避免直接平均百分比；
  \item 说明空值、未完成事件和分母为零在管理解释中的含义；
  \item 使用窗口函数计算排序、累计值、移动窗口和对象历史行为；
  \item 设计包含名称、业务定义、对象范围、分子分母、时间窗口、单位、责任人与 SQL 片段的指标字典；
  \item 写出可以被复核、被测试、被审计的指标查询。
\end{enumerate}

\section{从数据模型到指标层}

SQL 常被误解为一种“取数工具”。在本课程中，SQL 的意义更深：它是把企业数据模型转化为共同事实基础的主要语言。看板上的销售额、经营分析报告中的延期率、预警规则中的高风险订单数，本质上都依赖某种可复核的查询表达。

第2讲强调事实表粒度。第3讲继续使用同一原则。若订单表的粒度是一行一个订单，那么订单数可以直接按行计数；若事实表的粒度是一行一个订单明细项，那么订单数需要按订单编号去重；若同一张表混入订单头、订单行和库存快照，就会导致重复计数。SQL 语法本身不会阻止这种错误，只有明确的业务粒度和指标定义才能防止它。

\begin{definition}[指标层]
指标层是组织对管理指标的统一表达层。它规定指标名称、业务定义、对象范围、事实粒度、分子分母、维度、时间窗口、单位、刷新频率、权限和责任人，并提供可复用的 SQL 或语义配置。
\end{definition}

本讲把 SQL 放在如下链条中理解：
\[
\text{业务对象}
\rightarrow
\text{事实粒度}
\rightarrow
\text{查询表达}
\rightarrow
\text{指标口径}
\rightarrow
\text{管理解释}
\rightarrow
\text{行动审计}.
\]

在管理场景中，一个 SQL 结果是否“正确”，不仅取决于语法是否能运行，还取决于它是否回答了正确的业务问题，是否使用稳定的数据范围，是否保留了分子分母，是否允许他人复核。

\section{SQL 查询的关系语义}

关系模型把数据看作关系，即带有属性名的元组集合。SQL 查询可以理解为在关系上的变换。尽管实际数据库会优化执行计划，但从逻辑上看，一个典型查询依次经历表来源、连接、过滤、分组、聚合、筛选和输出。

\begin{definition}[SQL 查询的逻辑处理顺序]
一个常见 SQL 查询的逻辑顺序可概括为：
\[
\texttt{FROM/JOIN}
\rightarrow
\texttt{WHERE}
\rightarrow
\texttt{GROUP BY}
\rightarrow
\texttt{HAVING}
\rightarrow
\texttt{SELECT}
\rightarrow
\texttt{ORDER BY}.
\]
这一顺序表示查询语义的理解方式，不一定等于数据库系统实际执行计划。
\end{definition}

例如，下面的查询按区域计算 2026 年 3 月订单收入：

\begin{lstlisting}[caption={区域销售额查询}]
SELECT
  c.region,
  SUM(o.quantity * p.unit_price) AS revenue
FROM orders AS o
JOIN customers AS c
  ON o.customer_id = c.customer_id
JOIN products AS p
  ON o.product_id = p.product_id
WHERE o.order_date >= DATE '2026-03-01'
  AND o.order_date <  DATE '2026-04-01'
GROUP BY c.region
ORDER BY revenue DESC;
\end{lstlisting}

这段查询至少包含四个管理含义。第一，分析对象是 2026 年 3 月订单；第二，订单通过客户维度获得地区，通过产品维度获得售价；第三，收入按地区汇总；第四，输出排序用于比较区域经营表现。

\begin{remark}
\texttt{WHERE} 在聚合前过滤原始事实记录，\texttt{HAVING} 在聚合后过滤分组结果。若把二者混淆，就可能把“筛选订单”和“筛选地区”混为一谈。前者改变事实集合，后者改变展示集合。
\end{remark}

\section{关系代数视角}

SQL 之所以适合作为指标层语言，是因为它建立在关系模型和集合运算之上。虽然学生在实践中主要书写 SQL 语句，但理解其关系代数含义有助于判断查询是否改变了事实集合、是否丢失对象、是否重复计数。

\begin{table}[htbp]
\centering
\caption{关系代数操作与 SQL 子句}
\begin{tabular}{p{0.20\linewidth}p{0.29\linewidth}p{0.36\linewidth}}
\toprule
关系操作 & SQL 表达 & 管理含义\\
\midrule
选择 \(\sigma\) & \texttt{WHERE} & 限定事实集合，例如只保留某月订单。\\
投影 \(\pi\) & \texttt{SELECT} & 选择输出属性，例如地区和收入。\\
连接 \(\Join\) & \texttt{JOIN} & 补充业务对象属性，例如客户地区和产品价格。\\
分组聚合 \(\gamma\) & \texttt{GROUP BY} & 将事实提升到分析粒度，例如区域级 KPI。\\
排序 \(\tau\) & \texttt{ORDER BY} & 为比较和展示建立顺序。\\
\bottomrule
\end{tabular}
\end{table}

前一节的区域收入查询可以抽象为
\[
\gamma_{\text{region};\, \sum(q\cdot price)\rightarrow revenue}
\left(
\sigma_{\text{March 2026}}
(\text{orders}\Join \text{customers}\Join \text{products})
\right).
\]
这个表达式提醒我们，指标计算并不是随意拼接字段，而是在明确事实集合后做连接、选择和聚合。若事实集合错了，后续聚合再精细也无法恢复正确含义。

\begin{proposition}[声明式查询与可审计性]
SQL 的声明式特征使查询可以被复核：读者能够从查询中识别事实来源、过滤条件、连接键、分组维度和聚合方式。因此，清晰的 SQL 本身就是指标审计的一部分。
\end{proposition}

\section{事实粒度与分析粒度}

指标计算首先要声明粒度。粒度不清，是 KPI 错误的根源之一。

\begin{definition}[事务粒度、分析粒度与展示粒度]
事务粒度是源事实表中一行记录代表的业务事实；分析粒度是 SQL 查询分组后每一行结果代表的业务层级；展示粒度是看板、报告或管理会议中最终呈现的层级。
\end{definition}

例如，订单明细事实表的事务粒度可能是一行一个订单明细项；按地区、月份聚合后的分析粒度是一行一个地区月份；在管理层月报中展示的粒度可能是一行一个区域。三者可以不同，但必须被明确说明。

\begin{table}[htbp]
\centering
\caption{三类粒度的例子}
\begin{tabular}{p{0.21\linewidth}p{0.31\linewidth}p{0.34\linewidth}}
\toprule
粒度类型 & 示例 & 典型风险\\
\midrule
事务粒度 & 一行一个订单明细项。 & 误把明细行数当订单数。\\
分析粒度 & 一行一个地区、月份。 & 聚合维度遗漏导致混合口径。\\
展示粒度 & 一行一个区域或业务线。 & 从低层百分比直接平均得到高层百分比。\\
\bottomrule
\end{tabular}
\end{table}

SQL 查询写作的第一步不是输入 \texttt{SELECT}，而是写下如下句子：“本查询以什么事实表为基础？事实表一行代表什么？最终输出一行代表什么？”这句话比任何语法细节都更接近管理指标的核心。

\section{连接、粒度与重复计数}

连接是 SQL 中最容易产生管理分析错误的操作。内连接保留两张表中能够匹配的记录，左连接保留左表全部记录并补充右表信息，外连接可用于识别匹配与未匹配的边界。语法本身不难，真正困难的是判断连接是否改变了事实表粒度。

假设订单表 \(\texttt{orders}\) 的粒度是一行一个订单，产品表 \(\texttt{products}\) 的粒度是一行一个产品。如果 \(\texttt{products.product\_id}\) 唯一，则
\[
|\texttt{orders}\Join\texttt{products}|=|\texttt{orders}|.
\]
这说明连接只是补充产品属性，没有复制订单事实。若产品表中同一个 \(\texttt{product\_id}\) 出现两次，则连接后订单行数会膨胀，销售额、订单数和延期率都会被放大或扭曲。

\begin{longtable}{p{0.21\linewidth}p{0.33\linewidth}p{0.34\linewidth}}
\caption{连接错误与指标后果}\\
\toprule
问题 & 数据原因 & 指标后果\\
\midrule
\endfirsthead
\toprule
问题 & 数据原因 & 指标后果\\
\midrule
\endhead
事实行膨胀 & 维度键不唯一，或多对多关系未预聚合。 & 销售额、订单量、客户数被重复计算。\\
事实行丢失 & 使用内连接，但维度缺失或外键无法匹配。 & 低估业务规模，遗漏异常对象。\\
时间错配 & 使用当前维度属性解释历史事实。 & 历史地区、等级、价格口径被改写。\\
粒度混合 & 订单头、订单行、库存快照混在同一聚合层。 & KPI 分母不清，比较失真。\\
\bottomrule
\end{longtable}

\begin{proposition}[连接前后的行数检查]
若 \(F\) 是事实表，\(D\) 是维度表，且 \(D\) 在连接键上唯一并覆盖 \(F\) 的外键，则 \(F\) 与 \(D\) 的内连接不会改变事实行数。若连接后行数增加，通常说明维度键不唯一；若连接后行数减少，通常说明存在无法匹配的外键。
\end{proposition}

因此，在写任何管理指标 SQL 前，应先回答三个问题：事实表一行代表什么？连接后的行数是否仍代表同一事实？被连接的维度表在连接键上是否唯一？

\section{KPI 的形式化口径}

KPI 不是孤立数字，而是一组被治理的定义。一个管理指标可以形式化为
\[
M=(O,\mathcal G,\mathcal T,N,D,u,\rho),
\]
其中 \(O\) 是被度量的对象集合或事件集合，\(\mathcal G\) 是分组维度，\(\mathcal T\) 是时间窗口，\(N\) 是分子或可加度量，\(D\) 是分母，\(u\) 是单位，\(\rho\) 是业务责任或解释责任。

以准时交付率为例。对象集合 \(O\) 是已交付订单，时间窗口 \(\mathcal T\) 可以是订单日期或承诺交付日期所在月份，分子 \(N\) 是准时交付订单数，分母 \(D\) 是已交付订单数，单位 \(u\) 是百分比，责任 \(\rho\) 可能涉及供应链、仓储和物流团队。

\begin{table}[htbp]
\centering
\caption{样例 KPI 字典}
\begin{tabular}{p{0.18\linewidth}p{0.25\linewidth}p{0.39\linewidth}}
\toprule
指标 & 业务含义 & 计算口径\\
\midrule
销售额 & 订单带来的标准收入。 & \(\sum q_i \cdot price_i\)\\
毛利 & 收入扣除标准成本后的贡献。 & \(\sum q_i(price_i-cost_i)\)\\
毛利率 & 收入中转化为毛利的比例。 & \(\sum GP_i / \sum Revenue_i\)\\
准时交付率 & 已交付订单中按承诺日期交付的比例。 & \(\sum I(a_i\le p_i)/n\)\\
平均延期天数 & 订单超过承诺日期的平均延迟。 & \(\sum \max(a_i-p_i,0)/n\)\\
\bottomrule
\end{tabular}
\end{table}

注意，毛利率和准时交付率是比率型指标。比率型指标通常不可加，不能直接把不同区域、产品或月份的百分比再取简单平均。正确做法是保留分子和分母，在目标粒度上重新计算：
\[
\frac{\sum_{g}N_g}{\sum_{g}D_g}
\ne
\frac{1}{|\mathcal G|}\sum_g\frac{N_g}{D_g}.
\]

\begin{remark}
左边是按事实量加权后的总体比率，右边是分组百分比的平均数。除非每个分组分母相同，否则二者通常不同。
\end{remark}

\section{条件聚合与安全除法}

SQL 中的 KPI 通常通过条件聚合表达。下面的查询计算区域级销售额、毛利、准时交付率与平均延期天数。

\begin{lstlisting}[caption={区域级运营 KPI 查询}]
WITH order_enriched AS (
  SELECT
    o.order_id,
    c.region,
    o.quantity * p.unit_price AS revenue,
    o.quantity * (p.unit_price - p.unit_cost) AS gross_profit,
    CASE
      WHEN o.actual_date <= o.promised_date THEN 1
      ELSE 0
    END AS on_time_flag,
    GREATEST(date_diff('day', o.promised_date, o.actual_date), 0) AS delay_days
  FROM orders AS o
  JOIN customers AS c
    ON o.customer_id = c.customer_id
  JOIN products AS p
    ON o.product_id = p.product_id
  WHERE o.actual_date IS NOT NULL
)
SELECT
  region,
  COUNT(*) AS delivered_orders,
  SUM(revenue) AS revenue,
  SUM(gross_profit) AS gross_profit,
  SUM(gross_profit) / NULLIF(SUM(revenue), 0) AS gross_margin,
  SUM(on_time_flag) * 1.0 / NULLIF(COUNT(*), 0) AS on_time_rate,
  AVG(delay_days) AS average_delay_days
FROM order_enriched
GROUP BY region
ORDER BY revenue DESC;
\end{lstlisting}

\texttt{NULLIF} 的作用是避免分母为零时产生错误或无意义结果。管理指标中，分母为零不是可以简单忽略的技术异常，而是需要解释的业务状态。例如某区域本月没有已交付订单，则准时交付率并不是 0，而是没有定义；报表应区分“没有发生业务”和“业务发生但没有准时”。

\begin{definition}[安全除法]
安全除法是指在计算比率指标时显式处理分母为零、分母为空或分子分母口径不一致的情况。常见做法是返回空值、单独标记或在指标字典中说明不可定义状态。
\end{definition}

\section{空值、未完成事件与业务状态}

SQL 中的 \texttt{NULL} 不等同于数值 0，也不等同于空字符串。它通常表示未知、缺失、尚未发生或不适用。在管理指标中，空值处理必须回到业务状态，而不能只按技术便利处理。

\begin{table}[htbp]
\centering
\caption{空值与业务解释}
\begin{tabular}{p{0.24\linewidth}p{0.28\linewidth}p{0.34\linewidth}}
\toprule
情形 & 常见 SQL 处理 & 管理解释\\
\midrule
未交付订单无实际日期 & 从准时率分母中排除，或单独形成未完成指标。 & 尚无交付结果，不应简单计为延期或准时。\\
区域没有已交付订单 & 比率返回 \texttt{NULL}。 & 指标不可定义，而不是表现为 0。\\
产品成本缺失 & 标记数据质量问题。 & 毛利无法可信计算，需要补数或剔除说明。\\
客户维度缺失 & 左连接保留事实并标为未知客户。 & 业务规模存在，但归属和分层解释受限。\\
\bottomrule
\end{tabular}
\end{table}

\begin{remark}
一个严谨的指标系统应区分三类状态：“业务没有发生”“业务发生但结果失败”“数据缺失导致无法判断”。这三者在管理上对应完全不同的行动。
\end{remark}

\section{窗口函数与经营节奏}

聚合查询会把多行压缩为一行。窗口函数则不同，它在保留行身份的同时，对与当前行相关的一组行进行计算。这个思想非常适合经营分析，因为管理者既需要看到每个订单，也需要知道它在客户历史、区域排名或时间趋势中的位置。

\begin{lstlisting}[caption={客户最近三笔订单收入}]
SELECT
  o.customer_id,
  o.order_id,
  o.order_date,
  o.quantity * p.unit_price AS revenue,
  SUM(o.quantity * p.unit_price) OVER (
    PARTITION BY o.customer_id
    ORDER BY o.order_date, o.order_id
    ROWS BETWEEN 2 PRECEDING AND CURRENT ROW
  ) AS customer_recent_revenue,
  ROW_NUMBER() OVER (
    PARTITION BY o.customer_id
    ORDER BY o.order_date, o.order_id
  ) AS customer_order_sequence
FROM orders AS o
JOIN products AS p
  ON o.product_id = p.product_id
ORDER BY o.customer_id, o.order_date, o.order_id;
\end{lstlisting}

窗口函数在管理分析中常用于：
\begin{itemize}
  \item 对客户、产品、供应商或区域进行排名；
  \item 计算累计销售额、累计延期订单数和累计缺货次数；
  \item 计算移动窗口收入、近三次订单金额或近七天投诉数；
  \item 识别同一对象的前后状态变化，例如订单状态流转或库存变化。
\end{itemize}

窗口函数的风险在于窗口定义不清。若 \texttt{ORDER BY} 没有唯一排序键，排名和前后项可能不稳定；若窗口框架没有明确说明行数或时间范围，累计值可能被误解为移动值；若分区维度错误，客户级行为会被混成全局行为。

\section{从 SQL 到可审计指标层}

在小型课堂实验中，学生可以把 SQL 写在脚本里。但在企业环境中，指标需要被治理。所谓可审计指标层，是指组织能够回答以下问题：这个指标由谁定义？它使用哪些源表？它的事实粒度是什么？它能按哪些维度下钻？它是否可以跨时间求和？它的 SQL 是否经过测试？当一个经营会议引用该指标时，其他团队能否复现同样结果？

一个课程级指标配置可以写成如下结构：

\begin{lstlisting}[caption={指标配置示例}]
metrics:
  - name: gross_margin
    label: Gross margin
    type: ratio
    numerator: gross_profit
    denominator: revenue
    grain: order
    time_dimension: order_date
    dimensions: [region, segment, category, channel]
    owner: finance
\end{lstlisting}

\begin{table}[htbp]
\centering
\caption{可审计指标层的最低字段}
\begin{tabular}{p{0.22\linewidth}p{0.60\linewidth}}
\toprule
字段 & 说明\\
\midrule
指标名称 & 组织中统一使用的名称。\\
业务定义 & 指标回答的管理问题。\\
事实粒度 & 计算依据的一行事实代表什么。\\
分子与分母 & 比率指标必须保存分子和分母定义。\\
时间窗口 & 使用订单日期、交付日期、快照日期还是入库日期。\\
维度范围 & 可以按哪些对象、组织层级和时间层级下钻。\\
责任人与版本 & 谁维护口径，何时修改，如何复核。\\
\bottomrule
\end{tabular}
\end{table}

指标层的核心价值并不只是减少重复代码。它让组织能够围绕同一事实讨论问题，并让模型、看板、报告和审计共用一套经过治理的定义。

\section{从指标结果到管理解释}

SQL 生成的是结果表，管理者需要的是解释和行动。一个指标结果进入经营会议前，至少应完成六步解释。

\begin{table}[htbp]
\centering
\caption{KPI 结果解释模板}
\begin{tabular}{p{0.20\linewidth}p{0.63\linewidth}}
\toprule
解释环节 & 需要回答的问题\\
\midrule
事实描述 & 指标在什么对象、时间窗口和维度上发生了什么变化？\\
比较基准 & 与上期、目标值、预算、同类区域或历史分位数相比如何？\\
口径说明 & 分子、分母、排除条件、空值和未完成事件如何处理？\\
异常定位 & 哪些客户、产品、渠道、供应商或流程节点贡献最大？\\
行动建议 & 可以采取哪些行动，行动对象和约束是什么？\\
复核路径 & SQL、输入表版本、指标字典和责任人能否被追溯？\\
\bottomrule
\end{tabular}
\end{table}

例如，若南区准时交付率下降，合格的解释不应只写“南区表现较差”。更好的写法是：以已交付订单为分母，南区 2026 年 3 月准时交付率低于全体水平；下降主要来自高优先级订单和某一产品类别；未交付订单未进入分母；下一步应检查供应商交期、仓库出库时延和物流路线，并对高价值高风险订单建立预警。这样的解释才把 SQL 输出连接回管理行动。

\section{第3周实验案例}

本周样例数据位于 \texttt{data/sample/}。表结构如下：
\begin{itemize}
  \item \texttt{customers.csv}：客户维度，一行一个客户；
  \item \texttt{products.csv}：产品维度，一行一个产品，包含标准售价和单位成本；
  \item \texttt{orders.csv}：订单事实，一行一个订单；
  \item \texttt{inventory.csv}：库存快照，一行一个产品、仓库、日期。
\end{itemize}

课堂实验分为四步：
\begin{enumerate}
  \item 加载 CSV 数据并检查行数、主键唯一性和外键匹配；
  \item 用 SQL 连接订单、客户和产品三张表，构造订单级收入、毛利和交付状态；
  \item 按区域、客户分层、产品类别和渠道计算 KPI；
  \item 将 KPI 口径写入 \texttt{metrics.yml}，并用 Python 复核关键指标。
\end{enumerate}

学生需要特别注意，\texttt{orders.csv} 中的销售额不是原始字段，而是由数量和产品价格计算得到；毛利也不是原始字段，而是由数量、售价和成本共同确定。这种设计是刻意的，因为管理分析经常需要把事实表与维度表连接后才能得到指标度量。

\section{配套代码}

第三周 Python 复核代码位于：
\begin{center}
\texttt{src/bdm\_decision/cases/week03\_sql\_kpi.py}
\end{center}

核心函数包括：
\begin{itemize}
  \item \texttt{enrich\_orders()}：连接订单、客户和产品，生成订单级收入、毛利和交付状态；
  \item \texttt{summarize\_by()}：在指定分组维度上计算订单数、收入、毛利率、准时交付率和延期天数；
  \item \texttt{rolling\_customer\_revenue()}：复现窗口函数思想，计算客户最近三笔订单收入。
\end{itemize}

运行方式：
\begin{lstlisting}[language=bash,caption={运行第三周案例代码}]
PYTHONPATH=src python3 src/bdm_decision/cases/week03_sql_kpi.py
\end{lstlisting}

配套 SQL 位于：
\begin{center}
\texttt{src/bdm\_decision/warehouse/build\_kpi\_tables.sql}
\end{center}

该脚本以 DuckDB 风格读取 CSV，建立 \texttt{order\_enriched}、\texttt{kpi\_by\_region}、\texttt{kpi\_by\_month} 和 \texttt{customer\_order\_window} 等视图，便于学生把讲义中的 SQL 逻辑与可运行实验对应起来。

\section{可审计 SQL 的最低要求}

可审计 SQL 并不意味着复杂，而意味着别人可以理解、复现和质疑。课程要求学生在写 KPI SQL 时至少满足以下规范：
\begin{enumerate}
  \item 每个查询说明事实表粒度；
  \item 每个连接说明连接键及其唯一性假设；
  \item 每个比率指标保留分子和分母；
  \item 每个时间窗口说明采用订单日期、交付日期还是快照日期；
  \item 对空值、未交付订单和分母为零的情况给出处理规则；
  \item 查询输出字段命名稳定，能够进入后续 dashboard 或报告；
  \item 至少提供一个行数检查或 Python 复核结果。
\end{enumerate}

这些规范也为第4讲数据质量奠定基础。数据质量不是在指标算完之后才出现的问题；相反，重复主键、缺失外键、异常日期和字段空值会直接改变 SQL 的输出。因此，SQL/KPI 与数据质量不是先后无关的两个模块，而是同一条指标可信度链条上的两个环节。

\section*{习题}
\addcontentsline{toc}{section}{习题}

\begin{exercise}
使用 \path{data/sample/orders.csv}、\path{customers.csv} 和 \path{products.csv}，按客户分层计算销售额、毛利率和准时交付率。要求写出 SQL，并说明事实表粒度、分组粒度、分子、分母和时间窗口。
\end{exercise}

\begin{exercise}
解释为什么不能直接平均各区域毛利率来得到总体毛利率。请构造一个两区域的小例子，给出错误算法和正确算法的数值差异。
\end{exercise}

\begin{exercise}
写出一个 SQL 查询，找出每个客户最近三笔订单的累计收入。请说明 \texttt{PARTITION BY}、\texttt{ORDER BY} 和窗口框架分别承担什么作用。
\end{exercise}

\begin{exercise}
对 \texttt{orders} 与 \texttt{products} 的连接结果做行数检查。若连接后行数增加，说明可能出现了什么问题？若连接后行数减少，又可能是什么原因？
\end{exercise}

\begin{exercise}
为“高优先级订单延期率”设计一条指标字典记录，至少包含对象范围、事实粒度、分子、分母、时间窗口、单位、维度和责任部门。
\end{exercise}

\section*{本讲小结}
\addcontentsline{toc}{section}{本讲小结}

本讲建立了 SQL 与管理指标之间的联系。SQL 的核心不是机械语法，而是关系数据上的声明式表达；KPI 的核心不是单个数字，而是对象范围、粒度、分子分母、时间窗口、单位和责任的组合。连接决定事实是否被复制或遗漏，聚合决定指标落在哪个业务层级，空值处理决定指标能否被正确解释，窗口函数决定经营节奏如何被表达，指标层决定同一个组织是否能围绕同一口径讨论问题。

进入下一讲后，课程将继续追问：即使 SQL 口径正确，输入数据是否可信？主键是否重复？外键是否缺失？日期是否异常？指标口径是否被不同部门改写？这些问题构成数据质量与指标治理的核心。

\section*{常见误区}
\addcontentsline{toc}{section}{常见误区}

\begin{enumerate}
  \item 把 SQL 运行成功等同于指标正确。语法正确只是最低要求，指标还需要业务口径正确。
  \item 把百分比当作可加度量。比率型指标应在目标粒度重新聚合分子和分母。
  \item 在没有检查维度键唯一性的情况下连接维度表，导致事实行重复。
  \item 把未交付订单直接计入准时交付率分母，混淆“尚无结果”和“结果失败”。
  \item 把 \texttt{NULL} 当作 0，掩盖业务未发生、数据缺失或不适用之间的差异。
  \item 只保存最终 KPI 数字，不保存生成该数字的 SQL、时间窗口和输入数据版本。
  \item 在窗口函数中缺少稳定排序键，导致排名、前后项或移动窗口结果不可复现。
\end{enumerate}

\addcontentsline{toc}{chapter}{参考文献}
\begin{thebibliography}{99}
\bibitem{codd1970}
Codd, E. F. (1970).
\newblock A relational model of data for large shared data banks.
\newblock Communications of the ACM, 13(6), 377--387.

\bibitem{chamberlin1974}
Chamberlin, D. D., and Boyce, R. F. (1974).
\newblock SEQUEL: A structured English query language.
\newblock Proceedings of the 1974 ACM SIGFIDET Workshop on Data Description, Access and Control, 249--264.

\bibitem{date2003}
Date, C. J. (2003).
\newblock An Introduction to Database Systems.
\newblock Addison-Wesley.

\bibitem{kimball2013}
Kimball, R., and Ross, M. (2013).
\newblock The Data Warehouse Toolkit: The Definitive Guide to Dimensional Modeling.
\newblock Wiley.

\bibitem{iso9075}
International Organization for Standardization and International Electrotechnical Commission. (2016).
\newblock ISO/IEC 9075: Information technology -- Database languages -- SQL.

\bibitem{few2006}
Few, S. (2006).
\newblock Information Dashboard Design.
\newblock O'Reilly Media.

\bibitem{karwin2010}
Karwin, B. (2010).
\newblock SQL Antipatterns: Avoiding the Pitfalls of Database Programming.
\newblock Pragmatic Bookshelf.
\end{thebibliography}

\end{document}
